\par\Needspace{8\baselineskip}
\begin{center}[13.]\par\medskip\textit{Grundgesetz für alle Wechselwirkungen galvanischer Ströme.}\par\smallskip(Gefunden im Juli 1835.)\end{center}
Wenn sich zwei Elektrizitätselemente relativ zueinander bewegen, ziehen sie sich anders an oder stoßen sich anders ab, als wenn sie relativ zueinander ruhen.
\SourceLeaf{623}\[
\begin{gathered}
\begin{aligned}
e,x,y,z&\qquad\text{Element und Koordinaten}\\
e',x',y',z'&
\end{aligned}\\[4pt]
(x'-x)^2+(y'-y)^2+(z'-z)^2=rr
\end{gathered}
\]
Gegenseitige Wirkung (Abstoßung)
\[
=\frac{ee'}{rr}\left\{1+k\left(\begin{gathered}
\left(\frac{\mathrm d(x'-x)}{\mathrm dt}\right)^2+\left(\frac{\mathrm d(y'-y)}{\mathrm dt}\right)^2\\
+\left(\frac{\mathrm d(z'-z)}{\mathrm dt}\right)^2-\tfrac{3}{2}\left(\frac{\mathrm dr}{\mathrm dt}\right)^2
\end{gathered}\right)\right\}
\]
wobei $\BareRoot{\frac{1}{k}}$ eine bestimmte Geschwindigkeit bezeichnet.

\begin{center}[14.]\end{center}
Auf \textit{andere} Weise lässt sich das \textit{Grundgesetz} wie folgt darstellen.

Seien $P$ und $P'$ Punkte in zwei Strömen, $x,y,z$ und $x',y',z'$ ihre Koordinaten und $r$ ihr Abstand. Die Stromelemente $\mathrm ds,\mathrm ds'$ beginnen an diesen beiden Punkten.
\[
\begin{array}{rllll}
u&\text{Winkel}&\text{zwischen}&\mathrm ds\text{ und }\mathrm ds'\\
q&\text{-}&\text{-}&\mathrm ds\text{ und }PP'\\
q'&\text{-}&\text{-}&\mathrm ds'\text{ und }P'P
\end{array}
\]

Lässt man die Faktoren weg, die von der Stärke der Ströme abhängen, so üben die Elemente $\mathrm ds,\mathrm ds'$ eine gegenseitige Anziehung aus, die sich durch
\[
\frac{\mathrm ds\,.\,\mathrm ds'\,.\,(\cos u+\tfrac{3}{2}\cos q\,.\,\cos q')}{rr}
\]
messen lässt.

Setzen wir $\mathrm dx=\mathrm dy=0$, so ist diese Kraft
\[
=\frac{\mathrm ds\,.\,\mathrm dz'}{rr}-\tfrac{3}{2}\,.\,\frac{z'-z}{r^3}\mathrm ds\,.\,\mathrm ds'\,.\,\left(\frac{\mathrm dr}{\mathrm ds'}\right)
\]
oder, in Form der Kraftkomponenten, die parallel zu den Koordinatenachsen auf $\mathrm ds$ wirken,
\[
\begin{gathered}
\mathrm ds\,.\,\left\{\frac{(x'-x)\mathrm dz'}{r^3}-\frac{3(x'-x)(z'-z)\mathrm d'r}{2r^4}\right\}\\[4pt]
\mathrm ds\,.\,\left\{\frac{(y'-y)\mathrm dz'}{r^3}-\frac{3(y'-y)(z'-z)\mathrm d'r}{2r^4}\right\}\\[4pt]
\mathrm ds\,.\,\left\{\frac{(z'-z)\mathrm dz'}{r^3}-\frac{3(z'-z)(z'-z)\mathrm d'r}{2r^4}\right\}
\end{gathered}
\]
\SourceLeaf{624}oder, wenn man die Kräfte
\[
\begin{gathered}
-\mathrm ds\,.\,\mathrm d'\frac{(x'-x)(z'-z)}{2r^3}\\[4pt]
-\mathrm ds\,.\,\mathrm d'\frac{(y'-y)(z'-z)}{2r^3}\\[4pt]
-\mathrm ds\,.\,\mathrm d'\frac{(z'-z)(z'-z)}{2r^3}
\end{gathered}
\]
hinzufügt, was zulässig ist, sofern $\mathrm ds'$ ein Element eines geschlossenen Stroms ist, ergeben sich die Komponenten
\[
\begin{gathered}
\frac{\mathrm ds}{2r^3}((x'-x)\mathrm dz'-(z'-z)\mathrm dx')\\[4pt]
\frac{\mathrm ds}{2r^3}((y'-y)\mathrm dz'-(z'-z)\mathrm dy')\\[4pt]
\frac{\mathrm ds}{2r^3}((z'-z)\mathrm dz'-(z'-z)\mathrm dz')
\end{gathered}
\]
Die letzte ist offenbar $=0$. Diesen Kräften sind aber offenbar die folgenden gleichwertig, soweit die Wirkung auf $\mathrm ds$ betroffen ist.
\[
\begin{aligned}
\text{1. In der Richtung }PP'\; .\;.\;.\;.\;.\quad&\frac{\mathrm ds\,.\,\mathrm dz'}{2rr}=\frac{\mathrm ds\,.\,\mathrm ds'}{2rr}\cos u\\[4pt]
\text{2. in der Richtung parallel zu }\mathrm ds'\; .\;.\;.\;.\;.\quad&-\frac{\mathrm ds\,.\,\mathrm ds'}{2rr}\,.\,\cos q
\end{aligned}
\]
zu denen man noch hinzufügen darf
\[
\text{3. in der Richtung }\mathrm ds\; .\;.\;.\;.\;.\;.\quad+\frac{\mathrm ds\,.\,\mathrm ds'}{2rr}\,.\,\cos q'
\]
während auf $\mathrm ds'$ dann drei genau entgegengesetzte Kräfte wirken.

Seien $x,y,z$ die Koordinaten des wirkenden Stromelements und $0,0,0$ die Koordinaten des Elements, auf das eingewirkt wird. Werden Richtung und Stärke des ersten und des zweiten Elements nach den Koordinatenachsen durch $\xi',\eta',\zeta';\ \xi,\eta,\zeta$ angegeben, so beträgt die gesamte anziehende Kraft nach \textsc{Ampère}
\[
\frac{\xi\xi'+\eta\eta'+\zeta\zeta'}{rr}-\tfrac{3}{2}\frac{(x\xi+y\eta+z\zeta)(x\xi'+y\eta'+z\zeta')}{r^4}
\]
also ist die eine Kraftkomponente
\[
\begin{aligned}
&\frac{x}{r^3}(\xi'\mathrm dx+\eta'\mathrm dy+\zeta'\mathrm dz)-\frac{3x}{2r^4}(x\xi'+y\eta'+z\zeta')\mathrm dr\\[5pt]
&=\tfrac{1}{2}\xi'\mathrm d\frac{xx}{r^3}+\tfrac{1}{2}\eta'\mathrm d\frac{xy}{r^3}+\tfrac{1}{2}\zeta'\mathrm d\frac{xz}{r^3}\\
&\qquad+\frac{\eta'}{2r^3}(x\mathrm dy-y\mathrm dx)+\frac{\zeta'}{2r^3}(x\mathrm dz-z\mathrm dx)
\end{aligned}
\]
\SourceLeaf{625}oder, da man vollständige Differentiale weglassen kann,
\[
\frac{\eta'(x\eta-y\xi)+\zeta'(x\zeta-z\xi)}{2r^3}=\frac{x(\eta\eta'+\zeta\zeta')-y\xi\eta'-z\xi\zeta'}{2r^3}
\]
Hier darf man nun noch hinzufügen oder weglassen
\[
\frac{x\xi\xi'+y\eta\xi'+z\zeta\xi'}{2r^3}
\]
wodurch die Formel in Bezug auf beide Elemente symmetrisch wird. Wählen wir das Letztere, so erhalten wir die Kraft
\[
\frac{x(-\xi\xi'+\eta\eta'+\zeta\zeta')-y(\xi\eta'+\eta\xi')-z(\xi\zeta'+\zeta\xi')}{2r^3}
\]
Dies erklärt sich durch eine Kraft, die von den \textit{relativen} Bewegungen $\alpha,\beta,\gamma$ abhängt
\[
=-\frac{x(\alpha\alpha+\beta\beta+\gamma\gamma)-2\alpha(x\alpha+y\beta+z\gamma)}{2r^3}
\]

Hier kommen wir nun mit einer einzigen Zusatzkraft aus, die in Richtung $\mathrm{eb}$ mit der Stärke $\frac{vv}{kk}ee$ wirkt, wenn $\mathrm{ea}$ die Richtung der relativen Bewegung $v$ ist.
\[
\mathrm{aeb}=180^0-\mathrm{aeE},\qquad\mathrm{eE}\text{ positiv}
\]
\begin{center}\begin{tikzpicture}[x=.115mm,y=.115mm,line width=.65pt,>=stealth]
  \coordinate (E) at (0,0);
  \coordinate (e) at (551,0);
  \coordinate (back) at (434,-77);
  \coordinate (a) at (780,138);
  \coordinate (b) at (685,222);
  \draw (E)--(e);
  \draw[->] (back)--(e)--(a);
  \draw[dash pattern=on 3pt off 2.4pt,->] (e)--(b);
  \filldraw[fill=white] (E) circle[radius=2.1pt];
  \filldraw[fill=white] (e) circle[radius=2.1pt];
  \node[left=3pt,inner sep=0pt] at (E) {\bfseries E};
  \node[anchor=north west,yshift=-3pt,inner sep=0pt] at (e) {\bfseries e};
  \node[above=2pt,inner sep=0pt] at (a) {\bfseries a};
  \node[above=2pt,inner sep=0pt] at (b) {\bfseries b};
\end{tikzpicture}
\end{center}

Noch einfacher und ganz allgemein lässt sich das Gesetz wie folgt ausdrücken.
\begin{center}\begin{tikzpicture}[x=.115mm,y=.115mm,line width=.65pt,>=stealth]
  \coordinate (E) at (0,0);
  \coordinate (A) at (428,139);
  \coordinate (B) at (796,17);
  \coordinate (C) at (433,-124);
  \draw (E)--(B);
  \draw[->] (E)--(A);
  \draw[dash pattern=on 3pt off 2.4pt,->] (E)--(C);
  \filldraw[fill=white] (E) circle[radius=2.1pt];
  \node[anchor=south east,xshift=-3pt,inner sep=0pt] at (E) {\bfseries E};
  \node[above=2pt,inner sep=0pt] at (A) {\bfseries A};
  \node[right=4pt,inner sep=0pt] at (B) {\bfseries B};
  \node[below=2pt,inner sep=0pt] at (C) {\bfseries C};
\end{tikzpicture}
\end{center}

Wirkt auf ein elektrisches Element $\mathrm E$ durch ein anderes in Richtung $\mathrm{EA}$ die Kraft $p$, sofern beide relativ zueinander ruhen, so kommt bei einer relativen Bewegung, deren Richtung%
\SourceLeaf{626}%
\space in Bezug auf $\mathrm E$ die Gerade $\mathrm{EB}$ ist und deren Geschwindigkeit $=v$ beträgt, zu jener Kraft noch eine zweite hinzu, mit
\[
\text{Stärke}\quad=\frac{pvv}{kk},\qquad\text{Richtung}\quad=\mathrm{EC}
\]
wobei $\mathrm{EA,EB,EC}$ in einer einzigen Ebene liegen und $\mathrm{EB}$ genau zwischen $\mathrm{EA}$ und $\mathrm{EC}$ liegt.

